\begin{problem}[Quasi-compactness of principal opens]
Prove that every principal open $D(f)$ is quasi-compact.  Deduce that every affine spectrum
$\Spec A$ is quasi-compact.
\end{problem}

\paragraph{Why this problem matters.}
Quasi-compactness is one of the first global finiteness properties of affine geometry.  It
comes directly from finite algebraic expressions, not from metric compactness.

\paragraph{Useful ideas before starting.}
Apply the cover criterion of P3 to an arbitrary cover $D(f)\subseteq\bigcup D(g_\lambda)$.

\paragraph{Guided hints.}
If $f^N$ lies in the ideal generated by all $g_\lambda$, one expression for $f^N$ uses only
finitely many of them.

\begin{solution}
Suppose
\[
D(f)\subseteq\bigcup_{\lambda\in\Lambda}D(g_\lambda).
\]
By P3 there is $N\ge1$ with
\[
f^N\in(g_\lambda:\lambda\in\Lambda).
\]
By definition of an ideal generated by a family, there are finitely many indices
$\lambda_1,\dots,\lambda_r$ and coefficients $a_i\in A$ such that
\[
f^N=a_1g_{\lambda_1}+\cdots+a_rg_{\lambda_r}.
\]
Hence
\[
f\in\sqrt{(g_{\lambda_1},\dots,g_{\lambda_r})}.
\]
Another application of P3 gives the finite subcover
\[
D(f)\subseteq D(g_{\lambda_1})\cup\cdots\cup D(g_{\lambda_r}).
\]
Thus $D(f)$ is quasi-compact.  Taking $f=1$ yields
\[
\Spec A=D(1),
\]
so $\Spec A$ is quasi-compact.
\end{solution}

\paragraph{What happened mathematically.}
The finite-subcover property is exactly the statement that ideal membership is witnessed by
a finite sum.

\paragraph{Extensions.}
Prove that every quasi-compact open subset of $\Spec A$ is a finite union of principal opens.

\paragraph{Provenance.}
New canonical problem extending the legacy basis theorem to its fundamental finiteness
consequence.
